\section{Conclusion}

We audited \gittaskbench using only its official \processmetric and
\resultmetric outputs and found four consistent patterns across the evaluated
cohorts. Official success concentrates in a small subset of
domains---office-document, security, web, and speech---while image and video
remain dark. Models that tie on global totals succeed on different task
families, exposing specialization that a single score obscures. Provider models
are not interchangeable: they activate different families, and targeted
evaluation subsets can both overstate and miss activations that appear in
full-benchmark sweeps. The intermediate \processmetric$=$True,
\resultmetric$=$False regime clusters in structured-extraction families, mixing
threshold misses, parser failures, and evaluation-side errors rather than
behaving as random residue.

These findings lead to a practical recommendation. Repo-grounded agent
evaluations should report the official \processmetric/\resultmetric split and
family-level slices alongside any global pass rate. This requires no new
judges, taxonomies, or infrastructure---only preserving structure that the
benchmark already emits.

The recommendation is especially important as repo-grounded agent papers become
more incremental. When multiple systems cluster around small absolute gains,
readers need family-level slices to judge whether two systems improved in the
same way or simply activated different corners of the benchmark. Our findings
are demonstrated on \gittaskbench; validating whether similar structural
patterns hold on other repo-grounded benchmarks is a natural next step.
